```latex
\section*{Albondigas in Tomatensauce}
\ihead{}\chead{Personen: 4}\ohead{}
\ifoot{Vorbereitungszeit: 30 min}\cfoot{}\ofoot{Kochzeit: 30 min}

{\Large Zutaten}
\begin{multicols}{2}
\begin{itemize}
    \item \SI{500}{g} gemischtes Hackfleisch
    \item \SI{20}{g} gehackte Pinienkerne
    \item \num{2} Knoblauchzehen
    \item \num{1} Ei
    \item \SI{2-3}{EL} Paniermehl
    \item \SI{2}{EL} gehackte Petersilie
    \item Salz
    \item Pfeffer
    \item etwas Olivenöl zum Braten
    \item \SI{250}{g} Zwiebeln
    \item \num{3} Knoblauchzehen
    \item \num{2} TL getrockneter Thymian
    \item \num{2} Lorbeerblätter
    \item \num{1} TL Zucker
    \item \SI{800}{g} Tomaten im eigenen Saft
    \item \num{1} TL Gemüsebrühepulver
    \item Salz und Pfeffer
    \item \num{1} getrocknete Chilischote
\end{itemize}
\end{multicols}

\noindent
{\Large Zubereitung}\\
\\
\textit{Hackfleischbällchen:} Das Hackfleisch mit den Pinienkernen, Knoblauch, Ei, Paniermehl, Petersilie sowie Salz und Pfeffer ordentlich verkneten. Daraus kleine Bällchen formen. Es ergeben sich etwa 40 Stück, jeweils ungefähr so groß wie eine Walnuss.\\
Die Hackfleischbällchen in etwas Olivenöl bei mittlerer Hitze von allen Seiten anbraten, bis sie nicht mehr auseinanderfallen. Anschließend herausnehmen und beiseitestellen.\\
\textit{Sauce:} Die Zwiebeln fein würfeln und den Knoblauch in feine Scheiben schneiden. Beides im verbliebenen Bratfett glasig dünsten. Lorbeerblätter und Thymian dazugeben und weitere \SI{2}{min} dünsten.\\
Mit \SI{1}{TL} Zucker bestreuen und die grob zerkleinerten Tomaten dazugeben. Einmal aufkochen lassen. Anschließend mit Salz, Pfeffer, Gemüsebrühe und etwas Chili abschmecken.\\
Die Sauce etwa \SI{20}{min} bei niedriger Hitze köcheln lassen, bis sie sämig wird.\\
Zum Schluss die Hackfleischbällchen in die Sauce geben und darin gar ziehen lassen.\\
Die Albondigas warm oder heiß servieren. Dazu passen frisches Weißbrot und Aioli.
```
